% !TEX root = ../tfg_etsiinf_NombreApellidos.tex
\chapter{Impacto del trabajo} \label{chp:impacto}

Este capítulo analiza el impacto potencial del trabajo desde una perspectiva
técnica, académica y social. El análisis se centra en las posibilidades que
ofrece un entorno de aprendizaje por refuerzo reproducible para UAVs, así como
en sus limitaciones y riesgos cuando se traslada desde simulación hacia sistemas
robóticos reales.

\section{Impacto técnico y académico} \label{sct:impacto:tecnico}

El impacto más directo se produce en el ecosistema en el que el trabajo se
desarrolla. La infraestructura construida (entorno Gymnasium sobre Aerostack2,
reinicio certificado de episodios, guards de acción, supervisión de campañas
largas, vectorización validada y protocolo de evaluación reproducible) queda
publicada en repositorios abiertos y puede reutilizarse por el grupo de
investigación CVAR y por cualquier equipo que desee entrenar políticas de
control sobre Aerostack2 sin repetir el proceso de diagnóstico documentado en
esta memoria. Los hallazgos sobre el simulador, en particular la degradación
acumulativa asociada al servicio de reinicio y la persistencia indefinida de las
referencias de velocidad, constituyen además información accionable para los
mantenedores del proyecto original, ya que afectan a cualquier uso intensivo del
simulador más allá del aprendizaje por refuerzo.

En el plano metodológico, el trabajo aporta un ejemplo documentado de evaluación
homogénea entre un controlador aprendido y un controlador clásico: mismos
episodios generados por semillas, mismas métricas y una lectura explícita de
dónde gana cada enfoque. Varias decisiones se tomaron considerando este impacto:
la publicación del protocolo completo de evaluación, la conservación de los
experimentos fallidos como evidencia y la cuantificación de los fallos de
infraestructura en lugar de su omisión. Este material tiene también valor
docente, porque documenta una progresión completa en la que la mayor parte de la
dificultad no reside en el algoritmo de aprendizaje sino en la validez física
del entorno.

\section{Impacto social e industrial} \label{sct:impacto:social}

Los UAVs autónomos son una tecnología habilitadora en inspección de
infraestructuras, agricultura de precisión, logística y respuesta a emergencias.
Los controladores de navegación local como el entrenado en este trabajo son un
componente básico de esas aplicaciones, y la reducción de la barrera de entrada
para desarrollarlos sobre plataformas abiertas contribuye a que estas
capacidades no queden restringidas a soluciones propietarias.

Debe mantenerse, no obstante, cautela sobre el alcance real de lo demostrado.
Todos los resultados de este trabajo se obtienen en simulación, con una dinámica
simplificada y sin perturbaciones externas. La transferencia a vehículos reales
exige etapas adicionales de validación (robustez ante viento, ruido de sensores,
retardos de comunicación y fallos parciales) que quedan explícitamente fuera del
alcance. Presentar resultados de simulación como capacidades operativas sería
una extrapolación injustificada, y la memoria evita deliberadamente ese salto.


